% Section 5: Forecasting inside RL
\subsection{Forecasting inside RL}
\label{sec:forecasting_in_rl}

The previous sections used RL machinery to improve a forecaster. We now
invert the relation. The forecaster is the agent. Two flavours matter for
the chapter, namely trajectory generative models that treat the policy as
a forecast over future state-action sequences and counterfactual
forecasters that estimate the value of an alternative policy from logged
trajectories.

\subsubsection{Trajectory generative models}
\label{sec:trajectory_gen}

An offline reinforcement-learning dataset is a collection of trajectories
$\tau = (s_0, a_0, r_0, s_1, a_1, r_1, \dots, s_T, a_T, r_T)$ generated by
some behaviour policy. A trajectory generative model treats $\tau$ as a
sequence and trains a generative model on those sequences. The
\citet{Chen2021DT} Decision Transformer represents each trajectory as a
flat sequence
\begin{equation}
\tau = (\widehat R_0, s_0, a_0,\, \widehat R_1, s_1, a_1,\, \dots,\,
  \widehat R_T, s_T, a_T),
\label{eq:dt_sequence}
\end{equation}
where $\widehat R_t = \sum_{t' \ge t} r_{t'}$ is the return-to-go from
time $t$ onwards. A causally masked GPT-style transformer is trained to
predict the next action token given the last $K$ context tokens. At test
time, the operator specifies a desired return $R^\star$ and primes the
model with $\widehat R_0 = R^\star$ and the current state $s_0$. The model
generates the corresponding action, the environment returns a reward and
the next state, the operator decrements $R^\star$ by the realised reward,
and the loop repeats. No value function, no policy gradient, no critic
appears in training. The forecaster is the policy.

The empirical case rests on D4RL and Atari. On the $1\%$ DQN-replay slice
of Atari, the Decision Transformer attains gamer-normalised scores of
$267.5$ on Breakout against the $211.1$ of Conservative Q-Learning, the
incumbent offline-RL baseline. On the D4RL locomotion benchmark, the
Decision Transformer averages $74.7$ across nine task-dataset
combinations against CQL's $54.2$ and behavioural cloning's
$47.7$.\footnote{The same trajectory-as-sequence reframing was independently
proposed by \citet{Janner2021TT}, with beam search over tokenised states
and actions replacing return conditioning. Subsequent diffusion variants
\citep{Janner2022diffuser,Ajay2023} replace autoregressive generation with
iterative denoising of an entire trajectory and dispense with dynamic
programming entirely. The methodological move is the same. Plan by
sampling from a generative model of structured objects rather than by
maximising a learned value function.} For economists familiar with
structural estimation, the conceptual draw is that a forecast distribution
over trajectories doubles as the policy itself. The conditioning variable
plays the role of a behavioural target rather than a payoff.

\subsubsection{Counterfactual forecasting for off-policy evaluation}
\label{sec:counterfactual_ope}

Off-policy evaluation asks the question every applied economist asks at
some point. Given a logged dataset generated by a behaviour policy
$\pi_{\mathrm{obs}}$, what would the return have been under an alternative
policy $\widetilde\pi$? Importance sampling estimators are unbiased under
common support but high variance. Model-based estimators are low variance
but biased by model misspecification. The counterfactual route of
\citet{Buesing2019} sits between the two.

The construction casts the partially observable Markov decision process as
a structural causal model. The transition mechanism is
$s_{t+1} = f^s(s_t, a_t, u^s_t)$ with exogenous noise $u^s_t$, the
observation mechanism is $o_t = f^o(s_t, u^o_t)$, and the action
mechanism is $a_t = f^\pi(h_t, u^a_t)$ given the history. Given an
observed trajectory $\hat\tau$, the counterfactual rule has three steps.

\begin{algorithm}[h]
\caption{Counterfactual off-policy evaluation \citep{Buesing2019}}
\label{alg:buesing_cfope}
\begin{algorithmic}[1]
\State \textbf{Abduction.} Infer the posterior $P(U \mid \hat\tau)$ over the
exogenous noise from the observed trajectory.
\State \textbf{Action.} Replace the behaviour-policy mechanism $f^\pi$ with
the target-policy mechanism $f^{\widetilde\pi}$ to form the modified SCM
$\mathcal M^{\mathrm{do}(\widetilde\pi)}_{\hat\tau}$ with the same posterior
on $U$.
\State \textbf{Prediction.} Sample $U \sim P(U \mid \hat\tau)$ and roll out
the modified SCM to obtain counterfactual trajectories $\tau^{\mathrm{cf}}$
and their returns.
\end{algorithmic}
\end{algorithm}

Lemma~1 of \citet{Buesing2019} establishes that the marginal of the
counterfactual distribution over $\hat\tau$ equals the interventional
distribution. The counterfactual estimator is unbiased under correct SCM
specification. On the partially-observable grid-world experiment of the
paper, the counterfactual estimator's evaluation error decreases
monotonically with the amount of off-policy data, while the model-based
estimator's error remains stuck at the bias of the dynamics model.

The categorical-action case raises an identification problem.
\citet{Oberst2019} show that multiple structural causal models can be
consistent with the same interventional distribution yet imply different
counterfactual outcomes, and resolve the ambiguity by assuming a
Gumbel-max sampling mechanism that satisfies a counterfactual-stability
condition.\footnote{The Pearl monotonicity condition resolves the binary
case. The Gumbel-max SCM is the natural extension to $k$-way categorical
outcomes. \citet{TangWiens2023} combine SCM-based counterfactual outcomes
with importance sampling to obtain variance reductions while preserving
unbiasedness under correct human annotation, with a sepsis-simulator
demonstration that delivers an RMSE of $0.013$ against $0.113$ for vanilla
per-decision IS.} The shared identifying assumption across this strand is
that the structural causal model is correctly specified. The Lucas
critique applies in its sharpest form when the agent's policy itself
enters the data-generating process. Identification by instruments or by
exogeneity restrictions, in the spirit of \citet{Chernozhukov2018}, is the
econometric remedy and remains an open research question for the
counterfactual-RL programme.
